% Generated draft: requires mathematical and source review.
\documentclass[11pt]{article}
\usepackage[margin=25mm]{geometry}
\usepackage{fontspec}
\setmainfont{Latin Modern Roman}
\usepackage{amsmath,amssymb,mathtools,unicode-math}
\setmathfont{Latin Modern Math}
\usepackage{enumitem,needspace,xurl,hyperref}
\setlength{\parindent}{0pt}
\setlength{\parskip}{5pt}
\setlength{\emergencystretch}{3em}
\newcommand{\N}{\mathbb{N}}
\newcommand{\Z}{\mathbb{Z}}
\newcommand{\Q}{\mathbb{Q}}
\newcommand{\R}{\mathbb{R}}
\newcommand{\C}{\mathbb{C}}
\newcommand{\F}{\mathbb{F}}
\newcommand{\A}{\mathbb{A}}
\newcommand{\PP}{\mathbb{P}}
\newcommand{\E}{\mathbb{E}}
\newcommand{\eps}{\varepsilon}
\newcommand{\dd}{\,\mathrm{d}}
\begin{document}
\section*{UnsolvedMath ID: 30004113}
\textbf{Source problem code:} catalogue entry
\subsection{Definitions and mathematical statement}
The supplied record asks an extent-of-determination question rather than stating a quantified conjecture. Its mathematical target is:
\begin{quote}
How far does Jantzen's sum formula determine the characters of simple highest-weight modules for reductive groups?
\end{quote}
For terminology only, suppose a setting of finite-dimensional rational representations of a reductive algebraic group \(G\) has been specified, together with a maximal torus \(T\) and a choice of positive roots. Write \(X(T)\) for the character lattice of \(T\). For such a representation \(M\), its formal character is
\[
\operatorname{ch} M=\sum_{\mu\in X(T)}\dim(M_\mu)e^\mu\in\Z[X(T)],
\]
where \(M_\mu\) is its \(\mu\)-weight space and \(e^\mu\) denotes the corresponding basis element of the group ring. Write \(L(\lambda)\) for the simple highest-weight module of dominant highest weight \(\lambda\), whenever this notation applies in the chosen setting. The requested determination concerns \(\operatorname{ch}L(\lambda)\) using information provided by Jantzen's sum formula.

These conventions explain a customary interpretation; the source record does not specify the ground field, its characteristic, the precise class of reductive groups or representations, the allowed highest weights, or the version of Jantzen's sum formula. It also does not define which auxiliary information may be used or what constitutes determination by the formula. Consequently, neither a universal uniqueness assertion nor a classification of the cases of determination can be extracted as a precise conjecture from this record alone. The question above is retained with these unresolved qualifications.
\subsection{Short English statement}
Determine the extent to which Jantzen's sum formula can recover the formal characters of simple highest-weight modules for reductive groups. The source leaves the representation-theoretic setting and the meaning of determination unspecified, so this remains a qualified research question requiring human mathematical review.
\subsection{Sources}
\begin{itemize}
\item Local UnsolvedMath registry: supplied complete source record, permanent ID \(30004113\), title \emph{Character Determination by Jantzen's Sum Formula}, exact source problem code \texttt{catalogue entry}, in the Oberwolfach Reports Open Problems collection. The recorded proposed year is 2019; no individual proposer is supplied.
\item Supplied source citation: \emph{Representations of Finite Groups (2020)}, Oberwolfach Reports, DOI 10.4171/owr/2019/14, pp.\ 841--895. Supplied source and literature link, labelled ``Oberwolfach representation theory report'': \url{https://doi.org/10.4171/owr/2019/14}. The supplied background describes these reports as records of mathematical workshops and their open-problem discussions. No problem-specific page or passage is supplied.
\item Literature assessment dated exactly \texttt{2026-08-22}: the supplied curation labels the problem \texttt{partially\_solved} and states that Jantzen's sum formula resolves the relevant questions in several low-rank types and reduces further cases to finite ambiguities, while the report records additional higher-rank complications and no general solution. This is supplied dated curation, not an independently verified literature assessment; the background explicitly calls for independent verification of time-sensitive status and cited claims.
\item The supplied statement status is \texttt{corrected\_verified}. Its accompanying verification note reports manual semantic, typographic, and LaTeX review and claims that the clean statement is self-contained and preserves the target without unsupported additions. These are supplied registry assertions, not independent verification of the formulation or its mathematical status.
\end{itemize}
\end{document}
